\section*{अध्याय \ref{ch:b2:precambrian-earth} --- प्रीकैम्ब्रियन पृथ्वी}

\begin{solution}{exo:b2:precambrian-earth:1}
(क) \omterm{def:b2:precambrian-earth:eons}{हेडियन} $4567 - 4031 = \qty{536}{Myr}$; \omterm{def:b2:precambrian-earth:eons}{आर्कियन} $4031 - 2500 = \qty{1531}{Myr}$;
\omterm{def:b2:precambrian-earth:eons}{प्रोटेरोज़ोइक} $2500 - 538.8 = \qty{1961}{Myr}$। (ख) $(4567 - 538.8)/4567 = 0.882$: लगभग सात
आठवाँ भाग। (ग) ज़िरकॉन के लिए $4567 - 4404 = \qty{163}{Myr}$; सबसे पुरानी चट्टानों के लिए
$4567 - 4030 = \qty{537}{Myr}$।
\end{solution}

\begin{solution}{exo:b2:precambrian-earth:2}
(क) गुणक $e^{\lambda t}$, $\lambda = \ln2/T$ और $t = \qty{3}{Gyr}$ के साथ, हैं 1.593
(\ce{^{238}U}), 19.18 (\ce{^{235}U}), 1.160 (\ce{^{232}Th}) और 5.292 (\ce{^{40}K}): योगदान
3.053, 1.573, 2.424 और \qty{4.329}{pW} थे, कुल \qty{11.38}{pW}, आज के
\qty{4.906}{pW} के विरुद्ध। (ख) $11.38/4.906 = 2.32$। (ग) आज थोरियम-232 अग्रणी है; तब
पोटैशियम-40 था।
\end{solution}

\begin{solution}{exo:b2:precambrian-earth:3}
(क) $1616 + 18.0\times(28 - 32.2) = \qty{1540}{\celsius}$। (ख) $32.2 + (1500 - 1616)/18.0 =
\qty{25.8}{\percent}$। (ग) संपूर्ण-शैल विश्लेषण किसी द्रव का संघटन नहीं है: किसी
\omterm{def:b2:precambrian-earth:greenstone}{कोमैटिआइट} प्रवाह में द्रव से ढेर हुए ओलिवीन क्रिस्टल हो सकते हैं, जो MgO जोड़ते हैं। उसका
\qty{38}{\percent} सबसे समृद्ध \omterm{def:b2:precambrian-earth:greenstone}{कोमैटिआइट} द्रवों के \qty{32}{\percent} से आगे है, और
उन प्रयोगों से भी आगे जिन पर रेखा टिकी है।
\end{solution}

\begin{solution}{exo:b2:precambrian-earth:4}
सबसे पुरानी चट्टानें, \qty{4.03}{Ga}; पहले स्ट्रोमैटोलाइट, \qty{3.43}{Ga}; पिलबारा के गुंबदों का अंत,
\qty{3.1}{Ga}; \omterm{def:b2:precambrian-earth:goe}{महा ऑक्सीकरण घटना}, लगभग 2.4 से \qty{2.1}{Ga}; अंतिम सामान्य \omterm{def:b2:precambrian-earth:greenstone}{कोमैटिआइट},
\qty{1.9}{Ga}; ऑक्सीजन की दूसरी वृद्धि, लगभग \qty{0.8}{Ga}; मैरिनोअन का अंत,
\qty{0.635}{Ga}; पहले एडियाकरन जीवाश्म, लगभग \qty{0.57}{Ga}।
\end{solution}

\begin{solution}{exo:b2:precambrian-earth:5}
(क) $\tfrac14\times1361\times0.706 = \qty{240.2}{W.m^{-2}}$ अवशोषित और $\sigma\times288^4 =
\qty{390.1}{W.m^{-2}}$, अतः $\varepsilon = 240.2/390.1 = 0.616$। (ख) $\tfrac14\times1361\times0.4
= \qty{136.1}{W.m^{-2}}$, $T^4 = 136.1/(0.616\times\num{5.670e-8})$, $T = \qty{249.8}{K}$।
(ग) \qty{275.5}{K} पर एल्बिडो $0.6 - 0.306\times12.5/20 = 0.4088$ है: $R = 340.25\times0.5912 =
\qty{201.2}{W.m^{-2}}$ और $E = 0.616\,\sigma\times275.5^4 = \qty{201.2}{W.m^{-2}}$। (घ) $R' =
340.25\times0.306/20 = \qty{5.21}{W.m^{-2}.K^{-1}}$, $E' = 4\varepsilon\sigma T^3 =
\qty{2.92}{W.m^{-2}.K^{-1}}$: $R' > E'$, अतः साम्य अस्थायी है
(\cref{prop:b2:precambrian-earth:ebm})।
\end{solution}

\begin{solution}{exo:b2:precambrian-earth:6}
(क) $1/(1 + 0.4\times2.4/4.567) = 0.826$; $1/(1 + 0.4\times0.65/4.567) = 0.946$। (ख) गर्म
अवस्था के लिए आज के सूर्य-प्रकाश का कम से कम 0.932 चाहिए: \qty{2.4}{Ga} पहले असंभव,
\qty{0.65}{Ga} पहले संभव। (ग) $0.826/0.946 = 0.873$: \qty{13}{\percent} मंद, उद्धृत
\qty{10}{\percent} के निकट, जो एक भिन्न सौर प्रतिरूप से आता है।
\end{solution}

\begin{solution}{exo:b2:precambrian-earth:7}
(क) स्टर्टियन $716.9 - 658.9 = \qty{58.0}{Myr}$; मैरिनोअन $644.5 - 635.4 = \qty{9.1}{Myr}$;
अंतर्हिमनदीय काल $658.9 - 644.5 = \qty{14.4}{Myr}$। (ख) $(58.0 + 9.1)/85 = 0.79$: काल का लगभग चार
पाँचवाँ भाग। (ग) उसका आरंभ केवल 649.9 और \qty{639.0}{Ma} के बीच सीमित है: इन
सीमाओं पर मैरिनोअन लगभग 3 और \qty{15}{Myr} के बीच चला।
\end{solution}

\begin{solution}{exo:b2:precambrian-earth:8}
(क) $\tau = \frac{0.05}{0.95}\times\frac{44.01}{28.97}\times\frac{\num{5.1e18}}{\num{0.26e12}} =
0.05263\times1.519\times\num{1.962e7} = \qty{1.57}{Myr}$। (ख) $0.05/\num{278e-6} = 180$ गुना।
(ग) दुगुना समय, \qty{3.14}{Myr}।
\end{solution}

\begin{solution}{exo:b2:precambrian-earth:9}
(क) पत्थर नरम पंक में गिरे और उसकी परतों को नीचे मोड़ दिया: वे पातित पत्थर हैं, जिन्हें
तैरती बर्फ़ ने छोड़ा। (ख) भूमध्य रेखा से \ang{10} पर समुद्र तक पहुँचती बर्फ़ किसी
हिमगोल का अपेक्षित चिह्न है, और यदि केवल ध्रुवीय क्षेत्र जमे हों तो उसे पाना बहुत कठिन है। (ग) एक \omterm{def:b2:precambrian-earth:snowball}{टोपी कार्बोनेट}:
यह हिमनदन के अचानक अंत को अंकित करता है, तीव्र अपक्षय और पिघले जल द्वारा शेल्फ़ों के
डूबने के साथ।
\end{solution}

\begin{solution}{exo:b2:precambrian-earth:10}
(क) $0.94 > 0.932$, गर्म क्रांतिक बिंदु: गर्म साम्य, $288\times0.94^{1/4} =
\qty{283.6}{K}$, \qty{283}{K} से ठीक ऊपर है; और $0.94 < 1.227$, अतः हिमगोल भी मौजूद है।
(ख) $249.8\times0.94^{1/4} = \qty{246.0}{K}$। (ग) $\varepsilon = 0.94\times1361\times0.4/(4\sigma
\times263^4) = 511.7/1085.1 = 0.472$। (घ) $\sigma\times263^4 = \qty{271.3}{W.m^{-2}}$, और $(0.616
- 0.472)\times271.3 = \qty{39}{W.m^{-2}}$।
\end{solution}

\begin{solution}{exo:b2:precambrian-earth:11}
(क) $S_{\text{m}}/S_{\text{f}} = (T_1/T_2)^4\,(1 - a_{\text{w}})/(1 - a_{\text{i}})$:
उत्सर्जकता कट जाती है। $(263/283)^4\times0.706/0.4 = 0.7459\times1.765 = 1.317$। (ख) बर्फ़ और खुले जल के बीच
एल्बिडो का छोटा अंतर, जैसे अधिक गहरे रंग की, धूल-भरी बर्फ़, या तापमान का चौड़ा
ढलान। (ग) हिमगोल इतना सूर्य-प्रकाश परावर्तित करता है कि वह हर उस सूर्य-प्रकाश या
ग्रीनहाउस के नीचे ठंडा रहता है जिस पर गर्म पृथ्वी जम जाती; उसे पिघलाने के लिए बल को
जमने वाले बल से पाश की पूरी चौड़ाई जितना अधिक होना चाहिए।
\end{solution}

\begin{solution}{exo:b2:precambrian-earth:12}
$T^*$ के निकट, $C\dot\delta = -(E' - R')\delta$, अतः $\delta \propto e^{-t/\tau}$, जहाँ $\tau =
C/(E' - R')$। गर्म साम्य पर $R' = 0$ और $E' = 4\times0.616\times\num{5.670e-8}\times
288^3 = \qty{3.34}{W.m^{-2}.K^{-1}}$; $C = 1025\times3992\times50 = \qty{2.05e8}{J.m^{-2}.K^{-1}}$,
अतः $\tau = \qty{6.1e7}{s}$, लगभग \qty{1.9}{yr}। इस प्रतिरूप में गर्म साम्य $R' =
0$ बनाए रखता है जब तक वह \qty{283}{K} तक न पहुँचे, अतः $\tau$ क्रांतिक बिंदु तक लगभग दो वर्ष के निकट रहता है, फिर
अवस्था लुप्त हो जाती है। सुचारु रूप से बदलते एल्बिडो के साथ, साम्य के मोड़ के निकट आने पर $R'$ बढ़ता है,
जहाँ $E' = R'$: $\tau$ असीम रूप से बढ़ता है, एक धीमापन जो क्रांतिक
बिंदु की घोषणा करता है।
\end{solution}

\begin{solution}{pb:b2:precambrian-earth:1}
\textbf{1.} $\num{0.0203e-6}\times0.9929\times\num{95.13e-6} = \qty{1.92}{pW.kg^{-1}}$
(\ce{^{238}U}); $\num{0.0203e-6}\times0.00711\times\num{568.47e-6} = \qty{0.0820}{pW.kg^{-1}}$
(\ce{^{235}U}); $\num{0.0795e-6}\times\num{26.28e-6} = \qty{2.09}{pW.kg^{-1}}$ (\ce{^{232}Th});
$\num{240e-6}\times\num{1.20e-4}\times\num{28.47e-6} = \qty{0.820}{pW.kg^{-1}}$ (\ce{^{40}K})।
\textbf{2.} \qty{4.91}{pW.kg^{-1}}।
\textbf{3.} $\ln2/T$: 0.155, 0.985, 0.0495 और \qty{0.555}{Gyr^{-1}}।
\textbf{4.} गुणक 1.52, 14.3, 1.14 और 4.48: 2.91, 1.17, 2.39 और \qty{3.67}{pW.kg^{-1}}।
\textbf{5.} \qty{10.1}{pW.kg^{-1}}, आज का 2.07 गुना।
\textbf{6.} पोटैशियम-40। उसकी अर्धायु, \qty{1.25}{Gyr}, छोटी है: तब से दो से अधिक अर्धायु
बीत चुकी हैं, और उस समय के \ce{^{40}K} का तीन चौथाई से अधिक भाग क्षयित हो चुका है।
\textbf{7.} $1/(1 + 0.4\times2.7/4.567) = 0.809$; $0.809\times1361 = \qty{1101}{W.m^{-2}}$।
\textbf{8.} $288\times0.809^{1/4} = \qty{273}{K}$, \qty{283}{K} से नीचे: मान्य गर्म अवस्था नहीं;
पृथ्वी जम जाएगी।
\textbf{9.} $(1101\times0.4/(4\times0.616\,\sigma))^{1/4} = \qty{237}{K}$।
\textbf{10.} $0.616\times0.809 = 0.498$।
\textbf{11.} $0.004/\num{278e-6} = 14$ से $0.03/\num{278e-6} = 108$ गुना।
\textbf{12.} सिलिकेट-अपक्षय का तापस्थापी: अपक्षय ठंडी जलवायु में धीमा होता है, जिससे
ज्वालामुखीय कार्बन डाइऑक्साइड संचित होती है, और गर्म जलवायु में तेज़ होता है।
\textbf{13.} आज का 0.826, 0.932 से नीचे: कोई गर्म अवस्था नहीं; आज की वायु वाली पृथ्वी एक
हिमगोल होगी।
\textbf{14.} $\frac{0.1}{0.9}\times1.519\times\num{1.962e7} = \qty{3.31}{Myr}$।
\textbf{15.} $\frac{0.2}{0.8}\times1.519\times\num{1.962e7} = \qty{7.45}{Myr}$।
\textbf{16.} \omterm{def:b2:precambrian-earth:snowball}{टोपी कार्बोनेट}, हिमनदीय निक्षेपों के ऊपर तीक्ष्ण रूप से स्थित एक परतदार डोलोस्टोन,
जो गर्म परवर्ती अवस्था में तीव्र अपक्षय से पोषित और पिघले जल से डूबे समुद्र से
अवक्षेपित हुआ।
\textbf{17.} लौह(II) केवल ऑक्सीजन-रहित जल में घुला रहता है: लोहा बर्फ़ के नीचे
अनॉक्सी महासागर में संचित होता है और जहाँ ऑक्सीजन पहुँचती है, उदाहरण के लिए पिघले
जल के साथ, वहाँ लौह(III) के रूप में अवक्षेपित होता है।
\textbf{18.} $(1616 - 1450)/(32.2 - 23) = \qty{18.0}{\celsius}$ प्रति प्रतिशत।
\textbf{19.} $1616 + 18.0\times(30 - 32.2) = \qty{1576}{\celsius}$।
\textbf{20.} \qty{1616}{\celsius}, मापा गया बिंदु।
\textbf{21.} $1616 - 18.0\times20.2 = \qty{1252}{\celsius}$।
\textbf{22.} $1616 - 20\times20.2 = \qty{1212}{\celsius}$: \qty{40}{\celsius} कम, लावाओं के बीच के अंतर का
दसवाँ भाग; निष्कर्ष चयन पर निर्भर नहीं करता।
\textbf{23.} \omterm{def:b2:precambrian-earth:greenstone}{कोमैटिआइट} प्रावार-स्तंभों से आए, जो परिवेशी प्रावार से गर्म थे, और गहराई में गलन की बड़ी मात्राओं से;
किसी लावा का तापमान उसके स्रोत का विभव तापमान नहीं है। समग्र
प्रावार 200 से \qty{300}{K} गर्म था।
\textbf{24.} सबसे गर्म \omterm{def:b2:precambrian-earth:greenstone}{कोमैटिआइट} लावा आज के कटकों के किसी प्राथमिक बेसाल्ट से लगभग
\textbf{\qty{364}{\celsius}} अधिक गर्म था।
\end{solution}
